\section{AMC 10 10 2000}
        \begin{enumerate}
        	\item (\textit{AMC 10 10 2000}) In the year $2001$, the United States will host the International Mathematical Olympiad. Let $I, M,$ and $O$ be distinct positive integers such that the product $I\cdot M \cdot O=2001$. What is the largest possible value of the sum $I+M+O$?


 $\textbf{(A)}\ 23 \qquad\textbf{(B)}\ 55 \qquad\textbf{(C)}\ 99 \qquad\textbf{(D)}\ 111 \qquad\textbf{(E)}\ 671$


 
 \textbf{Answer: }E



\textbf{Solution 1}

The prime factorization of $2001$ is $3 \cdot 23 \cdot 29$.  

 The largest factor of $2001$ is $2001$.  The only way to write $2001$ as a product of three integers including $2001$ is $2001 \cdot 1 \cdot 1$, which is not allowed since $1$ is repeated.

 The second largest factor of $2001$ is $23 \cdot 29 = 667$.  The only way to write $2001$ as a product of three integers including $667$ is $667 \cdot 3 \cdot 1$, which is allowed.  In this case, the sum is $667 + 3 + 1 = 671$.

 The third largest factor of $2001$ is $3 \cdot 29 = 87$.  So, if $2001$ was written as a product of three integers without $2001$ or $667$, the largest of the integers would be at most $87$ and the sum must therefore be less than $671$.

 Therefore, the largest possible sum of the three factors is $\boxed{\textbf{(E)}\ 671}$.

 



\textbf{Solution 2 (Verifying the dumb Statement)}

First, we need to recognize that a number is going to be largest only if, of the $3$ \href{https://artofproblemsolving.com/wiki/index.php?title=Factor}{factors}, two of them are small. If we want to make sure that this is correct, we could test with a smaller number, like $30$. It becomes much more clear that this is true, and in this situation, the value of $I + M + O$ would be $18$. Now, we use this process on $2001$ to get $667 * 3 * 1$ as our $3$ factors.Hence, we have $667 + 3 + 1 = \boxed{\text{(E) 671.}}$

 



\textbf{Solution 3}

The sum is the highest if two \href{https://artofproblemsolving.com/wiki/index.php?title=Factor}{factors} are the lowest.

 So, $1 \cdot 3 \cdot 667 = 2001$ and $1+3+667=671 \Longrightarrow \boxed{\text{(E) 671}}$.

 



\textbf{Solution 4 (Answer Choices)}

We see since $2 + 0 + 0 + 1$ is divisible by $3$, we can eliminate all of the first $4$ answer choices because they are way too small and get $\boxed{\text{(E) 671}}$ as our final answer.

 


 


	\item (\textit{AMC 10 10 2000}) $2000({2000}^{2000}) =$


 $\textbf{(A)}\ {2000}^{2001} \qquad \textbf{(B)}\ {4000}^{2000} \qquad \textbf{(C)}\ {2000}^{4000} \qquad \textbf{(D)}\ {4,000,000}^{2000} \qquad\textbf{(E)}\ {2000}^{4,000,000}$


 
 \textbf{Answer: }A



\textbf{Solution 1}

Recall that $a^{b}\cdot a^{c} = a^{b+c}$ for any positive real numbers $a, b, c$.

 By rules of exponents, $2000\left(2000^{2000}\right) = \left(2000^{1}\right)\left(2000^{2000}\right) = 2000^{2000+1} = \boxed{\textbf{(A) } 2000^{2001}}$.

 Solution edited by armang32324, integralarefun, j314andrews

 



\textbf{Solution 2}

Since $2000^2000$ is the result of multiplying $2000$ copies of $2000$ together, $2000\left(2000^{2000}\right)$ must be the result of multiplying $2000 + 1 = 2001$ copies of $2000$ together, which results in $\boxed{\textbf{(A) } 2000^{2001}}$.

 - itsj, edited by j314andrews

 


	\item (\textit{AMC 10 10 2000}) Each day, Jenny ate $20\%$ of the jellybeans that were in her jar at the beginning of that day. At the end of the second day, $32$ remained. How many jellybeans were in the jar originally?


 $\textbf{(A)}\ 40 \qquad\textbf{(B)}\ 50 \qquad\textbf{(C)}\ 55 \qquad\textbf{(D)}\ 60 \qquad\textbf{(E)}\ 75$


 
 \textbf{Answer: }B



\textbf{Solution 1 (Algebra)}

Let $x$ be the number of jellybeans originally in the jar. Then after one day, there were $0.8x$ jellybeans left, while after two days, there were $0.8(0.8x)=0.64x$ jellybeans left.  So $0.64x = 32$ and thus $x = \boxed{\textbf{(B)}\ 50}$.

 -armang32324, edited by j314andrews

 



\textbf{Solution 2}

At the end of the day, $80\% = \frac{4}{5}$ of the jellybeans from the start of the day remained.  So after one day, there were $32 \div \frac{4}{5} = 40$ jellybeans left, while at the start, there were $40 \div \frac{4}{5} = \boxed{\textbf{(B)}\ 50}$ jellybeans left.

 



\textbf{Solution 3 (Answer Choices)}

Starting with each answer choice (A) through (E) results in $25.6$, $32$, $35.2$, $38.4$, and $48$ jellybeans left after two days, respectively.  So $\boxed{\textbf{(B)} \ 50}$ is correct.

 


	\item (\textit{AMC 10 10 2000}) Chandra pays an on-line service provider a fixed monthly fee plus an hourly charge for connect time. Her December bill was $\$12.48$, but in January her bill was $\$17.54$ because she used twice as much connect time as in December. What is the fixed monthly fee?


 $\textbf{(A) } \$2.53 \qquad\textbf{(B) } \$5.06 \qquad\textbf{(C) } \$6.24 \qquad\textbf{(D) } \$7.42 \qquad\textbf{(E) } \$8.77$


 
 \textbf{Answer: }D



\textbf{Solution 1}

Chandra spent $\$17.54 - \$12.48 = \$5.06$ more in January than December.  Since she used twice as much connect time in January as she did in December, she must have spent $\$5.06$ on connect time in December.  Therefore, the flat fee was $\$12.58 - \$5.06 = \boxed{\textbf{(D)}\ \$7.42}$

 



\textbf{Solution 2}

Let $x$ be the fixed fee, and $y$ be the amount Chandra spent on connect time in December.  Then:

 
 \[x+y=12.48\]
 \[x+2y=17.54\]

 Solving this system of equations yields $(x, y) = (7.42, 5.06)$, so the fixed fee is $\boxed{\textbf{(D)}\ \$7.42}$.

 


	\item (\textit{AMC 10 10 2000}) Points $M$ and $N$ are the midpoints of sides $PA$ and $PB$ of $\triangle PAB$. As $P$ moves along a line that is parallel to side $AB$, how many of the four quantities listed below change?


 (a) the length of the segment $MN$


 (b) the perimeter of $\triangle PAB$


 (c) the area of $\triangle PAB$


 (d) the area of trapezoid $ABNM$


 
                    \begin{figure}[H]
                        \centering
                        \includegraphics[width=0.5\textwidth]{Problem/AMC_10_10_2000_5_Problem_Figure_1.png}
                    \end{figure}
                    


 $\textbf{(A)}\ 0 \qquad\textbf{(B)}\ 1 \qquad\textbf{(C)}\ 2 \qquad\textbf{(D)}\ 3 \qquad\textbf{(E)}\ 4$


 
 \textbf{Answer: }B



\textbf{Solution}

Examine each quantity to determine if it changes or remains constant.

 (a) By the midsegment theorem, $MN=\frac{1}{2}AB$, so it remains constant.

 (b) If $P$ moves far to the left or right, the perimeter of $\triangle PAB$ increases without bound, so it changes.

 (c) The height of $\triangle PAB$ remains constant, as does its base $AB$.  Therefore, the area of $\triangle PAB$ remains constant.

 (d) The height of $ABNM$ remains constant, as do the lengths of its bases $AB$ and $MN$.  Therefore, the area of $ABNM$ remains constant.

 Thus $\boxed{\textbf{(B)}\ 1}$ of these quantities changes.

 


	\item (\textit{AMC 10 10 2000}) The Fibonacci sequence $1, 1, 2, 3, 5, 8, 13, 21, \ldots$ starts with two $1$s, and each term afterwards is the sum of its two predecessors. Which one of the ten digits is the  last to appear in the units position of a number in the Fibonacci sequence?


 $\textbf{(A)}\ 0 \qquad\textbf{(B)}\ 4 \qquad\textbf{(C)}\ 6 \qquad\textbf{(D)}\ 7\qquad\textbf{(E)}\ 9$


 
 \textbf{Answer: }C



\textbf{Solution}

The units digit of any term after the second term can be found by summing the units digits of the previous two terms and taking the units digit of that sum (i.e. taking the sum of the previous two terms modulo $10$).

 The sequence of units digits therefore starts with $1,1,2,3,5,8,3,1,4,5,9,4,3,7,0,7,7,4,1,5,6$, and $\boxed{\mathbf{(C)}\ 6}$ appears last.

 


	\item (\textit{AMC 10 10 2000}) In rectangle $ABCD$, $AD=1$, $P$ is on $\overline{AB}$, and $\overline{DB}$ and $\overline{DP}$ trisect $\angle ADC$.  What is the perimeter of $\triangle BDP$?


 
                    \begin{figure}[H]
                        \centering
                        \includegraphics[width=0.5\textwidth]{Problem/AMC_10_10_2000_7_Problem_Figure_1.png}
                    \end{figure}
                    


 $\textbf{(A)}\ 3+\frac{\sqrt{3}}{3} \qquad\textbf{(B)}\ 2+\frac{4\sqrt{3}}{3} \qquad\textbf{(C)}\ 2+2\sqrt{2} \qquad\textbf{(D)}\ \frac{3+3\sqrt{5}}{2} \qquad\textbf{(E)}\ 2+\frac{5\sqrt{3}}{3}$


 
 \textbf{Answer: }B



\textbf{Solution}


                    \begin{figure}[H]
                        \centering
                        \includegraphics[width=0.5\textwidth]{Solution/AMC_10_10_2000_7_Solution_Figure_1.png}
                    \end{figure}
                    

 $AD=1$.

 Since $\angle ADC$ is trisected, $\angle ADP= \angle PDB= \angle BDC=30^\circ$.

 Thus, $PD=\frac{2\sqrt{3}}{3}$

 $DB=2$

 $BP=\sqrt{3}-\frac{\sqrt{3}}{3}=\frac{2\sqrt{3}}{3}$.

 Adding, we get $\boxed{\textbf{(B) }  2+\frac{4\sqrt{3}}{3}}$.

 This is due to the fact that ABD and ABP are $30-60-90$ triangles

 -Amon26

 



\textbf{Solution 2}

After computing $\overline{BP} = \frac{2\sqrt{3}}{3},$ observe that triangle $\triangle BPD$ is isosceles with $\angle DPB = \angle BPD.$ Therefore, using $120 - 30 - 30$ triangle properties, we see that the perimeter is just $(2+ \sqrt{3}) \cdot \frac{2\sqrt{3}}{3} = \boxed{\textbf{(B) }  2+\frac{4\sqrt{3}}{3}}.$

 


	\item (\textit{AMC 10 10 2000}) At Olympic High School, $\frac{2}{5}$ of the freshmen and $\frac{4}{5}$ of the sophomores took the AMC 10. Given that the number of freshmen and sophomore contestants was the same, which of the following must be true? 


 $\textbf{(A) }$ There are five times as many sophomores as freshmen. 


 $\textbf{(B) }$ There are twice as many sophomores as freshmen.


 $\textbf{(C) }$ There are as many freshmen as sophomores.


 $\textbf{(D) }$ There are twice as many freshmen as sophomores. 


 $\textbf{(E) }$ There are five times as many freshmen as sophomores.


 
 \textbf{Answer: }D



\textbf{Solution}

Let $f$ be the number of freshman and $s$ be the number of sophomores.

 $\frac{2}{5}f=\frac{4}{5}s$

 $2f = 4s$

 $f=2s$

 There are twice as many freshmen as sophomores.$\boxed{\text{D}}$

 


	\item (\textit{AMC 10 10 2000}) If $|x-2|=p$, where $x<2$, then $x-p=$


 $\textbf{(A)}\ -2 \qquad\textbf{(B)}\ 2 \qquad\textbf{(C)}\ 2-2p \qquad\textbf{(D)}\ 2p-2 \qquad\textbf{(E)}\ |2p-2|$


 
 \textbf{Answer: }C



\textbf{Solution 1}

When $x < 2,$ $x-2$ is negative so $|x - 2| = 2-x = p$ and $x = 2-p$.

 Thus $x-p = (2-p)-p = \boxed{\mathbf{(C)}\ \ensuremath{2-2p}}$.

 



\textbf{Solution 2 (Answer Choices)}

Suppose $x = -10$.  Then $p = 12$, so $x - p = -22$.  Substituting $p = 12$ into choices (A) through (E) yields $-2$, $2$, $-22$, $22$, and $22$ respectively, so $\boxed{\mathbf{(C)}\ \ensuremath{2-2p}}$ is correct.

 -Monkey\_King, edited by j314andrews

 


	\item (\textit{AMC 10 10 2000}) The sides of a triangle with positive area have lengths $4$, $6$, and $x$. The sides of a second triangle with positive area have lengths $4$, $6$, and $y$. What is the smallest positive number that is \textbf{not} a possible value of $|x-y|$?  


 $\textbf{(A)}\ 2 \qquad\textbf{(B)}\ 4 \qquad\textbf{(C)}\ 6 \qquad\textbf{(D)}\ 8 \qquad\textbf{(E)}\ 10$


 
 \textbf{Answer: }D



\textbf{Solution}

By the triangle inequality, $x + 4 > 6$, so $x > 2$. Also, by the triangle inequality, $4 + 6 > x$, so $x < 10$.  Thus $2 < x < 10$, and by the same logic $2 < y < 10$.  

 So $x$ and $y$ are strictly between $2$ and $10$.  The length of this interval, which is also the smallest value which cannot be equal to $|x-y|$, is $\boxed{\textbf{(D)}\ 8}$.  For $0 \leq d < 8$, $|x-y|=d$ can be achieved by setting $(x, y) = \left(6+\frac{d}{2}, 6-\frac{d}{2}\right)$.

 -Jayeed Mahmud, edited by j314andrews

 


	\item (\textit{AMC 10 10 2000}) Two different prime numbers between $4$ and $18$ are chosen. When their sum is subtracted from their product, which of the following numbers could be obtained?


 $\textbf{(A)}\ 21 \qquad\textbf{(B)}\ 60 \qquad\textbf{(C)}\ 119 \qquad\textbf{(D)}\ 180 \qquad\textbf{(E)}\ 231$


 
 \textbf{Answer: }C



\textbf{Solution 1 (Meta-solving)}

Both prime numbers are odd, so their product is odd and their sum is even.  Therefore, the obtained number is odd, and (B) and (D) are incorrect. 

 The largest two prime numbers between $4$ and $18$ are $13$ and $17$, and the obtained number must be less than the product of the two primes.  So any number obtained is less than $13 \cdot 17 = 221$, so (E) is too large and therefore incorrect.

 If the prime numbers chosen are $a$ and $b$, then the required difference is $(a - 1)(b - 1) - 1$ so the minimum value is $4 \cdot 6 - 1 = 23$ so (A) is incorrect.

 This leaves $\boxed{\textbf{(C)}\ 119}$, which is correct as $11 \cdot 13 - (11+13) = 119$.

 



\textbf{Solution 2}

Let $p$ and $q$ be the two primes.  

 Then the obtained number is $pq-(p+q)=pq-p-q=pq-p-q+1-1=(p-1)(q-1)-1$ by Simon's favorite factoring trick. 

 Since $p$ and $q$ must be odd, $(p-1)(q-1)-1$ must be odd, so (A), (B), and (D) are incorrect.

 If $(p-1)(q-1)-1 = 231$, then $(p-1)(q-1) = 232 = 2^3 \cdot 29$.  So either $p-1$ or $q-1$ must be divisible by $29$, and thus either $p$ or $q$ is at least $30$, which is too large.  So (E) is incorrect.

 This leaves $\boxed{\textbf{(C)}\ 119}$, which is obtained if $(p, q) = (11, 13)$.

 


	\item (\textit{AMC 10 10 2000}) Figures $0$, $1$, $2$, and $3$ consist of $1$, $5$, $13$, and $25$ nonoverlapping unit squares, respectively. If the pattern were continued, how many nonoverlapping unit squares would there be in figure 100?


 Figures:

                    \begin{figure}[H]
                        \centering
                        \includegraphics[width=0.5\textwidth]{Problem/AMC_10_10_2000_12_Problem_Figure_1.png}
                    \end{figure}
                    


 

$\textbf{(A)}\ 10401 \qquad\textbf{(B)}\ 19801 \qquad\textbf{(C)}\ 20201 \qquad\textbf{(D)}\ 39801 \qquad\textbf{(E)}\ 40801$


 
 \textbf{Answer: }C



\textbf{Solution 1}


                    \begin{figure}[H]
                        \centering
                        
\includegraphics[width=0.5\textwidth]{Solution/AMC_10_10_2000_12_Solution_Figure_1.png}
                    \end{figure}
                    

 Shade in the figure so that the outermost squares are shaded and the rest of the grid follows a checkerboard pattern.  Notice that figure $n$ has $(n+1)^2$ shaded unit squares and $n^2$ unshaded unit squares for a total of $(n+1)^2 + n^2$ squares.  Therefore, figure $100$ has $101^2+100^2=\boxed{\textbf{(C)}\ 20201}$ unit squares.

 



\textbf{Solution 2}

Notice that figure $n$ consists of a central unit square and $n$ rings, where for $1 \leq k \leq n$, the $k\mathrm{th}$ ring contains $4k$ unit squares.  So figure $100$ contains $1 + 4 + 8 + 12 + ... + 400 = 1 + 4(1+2+...+100)=1+4\left(\frac{100\cdot101}{2}\right) = \boxed{\textbf{(C)}\ 20201}$ unit squares.

 



\textbf{Solution 3}

We can see that each figure $n$ has a central box and 4 columns of $n$ boxes on each side of each square. Therefore, at figure 100, there is a central box with 100 boxes on the top, right, left, and bottom. Knowing that each quarter of each figure has a pyramid structure, we know that for each quarter there are $\sum_{n=1}^{100} n = 5050$ squares. $4 \cdot 5050 = 20200$. Adding in the original center box we have $20200 + 1 = \boxed{\textbf{(C) }20201}$.

 



\textbf{Solution 4}

Let $a_n$ be the number of squares in figure $n$.  We can easily see that 
 \[a_0=4\cdot 0+1\]
 \[a_1=4\cdot 1+1\]
 \[a_2=4\cdot 3+1\]
 \[a_3=4\cdot 6+1.\]See that we multiply the number we are on to the next consecutive number. Note that in $a_n$, the number multiplied by the 4 is the $n$th triangular number.  Hence, $a_{100}=4\cdot \frac{100\cdot 101}{2}+1=\boxed{\textbf{(C) }20201}$.~ edited by mathlover66

 



\textbf{Solution 5}

Let $f_n$ denote the number of unit cubes in a figure. We have
 \[f_0=1\]
 \[f_1=5\]
 \[f_2=13\]
 \[f_3=25\]
 \[f_4=41\]
 \[...\]

 Computing the difference between the number of cubes in each figure yields
 \[4,8,12,16,...\]It is easy to notice that this is an arithmetic sequence, with the first term being $4$ and the difference being $4$. Let this sequence be $a_n$

 From $f_0$ to $f_{100}$, the sequence will have $100$ terms. Using the arithmetic sum formula yields

 
 \[S_{100}=\frac{100[2\cdot 4+(100-1)4]}{2}\]
 \[=50(2\cdot 4+99\cdot 4)\]
 \[=50(101\cdot 4)\]
 \[=200\cdot 101\]
 \[=20200\]

 So $f_{100}=1+20200=\boxed{\textbf{(C) }20201}$ unit cubes.

 



\textbf{Solution 6 (Newton's Forward Differences)}

We know that 1 and 5 differ by 4, 5 and 13 differ by 8, and 13 and 25 differ by 12. Hence the differences are 4, 8, and 12, resp. And the differences of the differences area all 4. So by Newton's Forward Difference Formula, we get the 100th figure is(because Figure 0 exists) $\dbinom{101-1}{0}+4\dbinom{101-1}{1}+4\dbinom{101-1}{2}=20201$ or $\textbf{(C)}$-vsamc

 


 



\textbf{Solution 7 (Newton’s little formula)}

Newton’s little formula states that $a_n = A \binom{n-1}{0} + B \binom{n-1}{1} + C \binom{n-1}{2} + \cdots + K \binom{n-1}{m}$ if first term is $a_1$ and $A =$ first difference, $B =$ second difference, and so on. Hence we apply the formula (because we start at term 0, term 100 is $a_{101}$): $a_{101} = 5 \binom{100}{0} + 8 \binom{100}{1} + 4 \binom{100}{2} = \boxed{20201} = \textbf{(C)}.$~Peelybonehead~clarification by LeonidasTheConquerer

 



\textbf{Solution 8 (geometrical intuition)}

By taking figure $n$, putting the centers of each square on a grid, and rotating the figure 45 degrees, we get the following shape (figure 3 is shown for reference):
                    \begin{figure}[H]
                        \centering
                        \includegraphics[width=0.5\textwidth]{Solution/AMC_10_10_2000_12_Solution_Figure_2.png}
                    \end{figure}
                    There are two lattices of points (shown in red): one of side length $n$ (blue) and one of $n+1$ (green). The rest follows like in solution 1.

 -\href{https://artofproblemsolving.com/wiki/index.php?title=User:Integralarefun}{Integralarefun} (\href{https://artofproblemsolving.com/wiki/index.php?title=User\_talk:Integralarefun}{talk}) 18:21, 29 July 2023 (EDT)

 



\textbf{Solution 9}

If you do not want to base off the previous number in the sequence, you can come up with a formula based on which figure you are looking for:

 $1+x(x-2+x)$

 
                \begin{itemize}
                    \item However, this formula says that in figure 1, there is 1 square and so forth. Therefore, if we are looking to find $f(100)$, we will have to find $f(101)$.
                \end{itemize}


                This is the formula because to get from 1 to 5, we have +4. To get 1 to 13, it is +12. To get 1 to 25, it is +24. Assuming that 5 squares is figure 2 rather than in figure one (sorry, it just made better sense in my brain), then in figure two, it is 2*2 to get +4 on top of one. In figure “3”, to get +12, you have 3*4. In figure 4, to get +24, you have 4*6. Now, you can notice that in figure two, it is 2* (the distance to 2- which is 0- added onto 2) which is +2*2. In figure 3, to get +12, it is 3 (the figure number) * (the distance from 3 to 2- which is 1, added back onto the figure number). In figure 4, to get +24, you can get 4*6 (the figure number * (the distance from 4 to two- which is 2- added back onto the figure number, 6). 

 So to get figure 100, we have to do f(101) because I said that figure 1 was 1 rather than figure 0. So 1+ 101(101-2+101) which is 20201.

 -Smartgrowth

 


	\item (\textit{AMC 10 10 2000}) There are 5 yellow pegs, 4 red pegs, 3 green pegs, 2 blue pegs, and 1 orange peg to be placed on a triangular peg board. In how many ways can the pegs be placed so that no (horizontal) row or (vertical) column contains two pegs of the same color?


 
                    \begin{figure}[H]
                        \centering
                        \includegraphics[width=0.5\textwidth]{Problem/AMC_10_10_2000_13_Problem_Figure_1.png}
                    \end{figure}
                    


 

$\textbf{(A)}\ 0 \qquad\textbf{(B)}\ 1 \qquad\textbf{(C)}\ 5!\cdot 4!\cdot 3!\cdot 2!\cdot 1!  \qquad\textbf{(D)}\ \frac{15!}{5!\cdot 4!\cdot 3!\cdot 2!\cdot 1!} \qquad\textbf{(E)}\ 15!$


 
 \textbf{Answer: }B



\textbf{Solution}

In each column there must be one orange peg. In particular, in the rightmost column, there is only one peg spot, therefore an orange peg must go there. 

 In the second column from the right, there are two spaces for pegs. One of them is in the same row as the corner peg, so there is only one remaining choice left for the orange peg in this column.

 By similar logic, we can fill in the orange pegs as shown:

 
                    \begin{figure}[H]
                        \centering
                        \includegraphics[width=0.5\textwidth]{Solution/AMC_10_10_2000_13_Solution_Figure_1.png}
                    \end{figure}
                    

 After this we can proceed to fill in the whole pegboard, so there is only $1$ arrangement of the pegs. The answer is $\boxed{\text{B}}$

 
                    \begin{figure}[H]
                        \centering
                        \includegraphics[width=0.5\textwidth]{Solution/AMC_10_10_2000_13_Solution_Figure_2.png}
                    \end{figure}
                    

 


	\item (\textit{AMC 10 10 2000}) Mrs. Walter gave an exam in a mathematics class of five students.  She entered the scores in random order into a spreadsheet, which recalculated the class average after each score was entered.  Mrs. Walter noticed that after each score was entered, the average was always an integer.  The scores (listed in ascending order) were $71$, $76$, $80$, $82$, and $91$.  What was the last score Mrs. Walter entered?


 $\textbf{(A)}\ 71 \qquad\textbf{(B)}\ 76 \qquad\textbf{(C)}\ 80 \qquad\textbf{(D)}\ 82 \qquad\textbf{(E)}\ 91$


 
 \textbf{Answer: }C



\textbf{Solution 1}

The total of the $5$ scores is $71+76+80+82+91=400$.

 Since the average of the first $4$ scores is an integer, the sum of the first $4$ scores is divisible by $4$.  Since $400$ is divisible by $4$, the last score is also divisible by $4$, so it is either $76$ or $80$.

 If the last score is $76$, the sum of the first $4$ scores is $400-76=324$.  Since the average of the first $3$ scores is an integer, the sum of the first $3$ scores is divisible by $3$.  Since $324$ is divisible by $3$, the fourth score is also divisible by $3$.  But this is impossible since none of the test scores are divisible by $3$.

 Therefore, the last score is $\boxed{\textbf{(C)}\ 80}$.  A possible order for the scores is $76,82,91,71,80$.

 



\textbf{Solution 2}

We know the first sum of the first three numbers must be divisible by $3,$ so we write out all $5$ numbers $\pmod{3}$, which gives $2,1,2,1,1,$ respectively. Clearly, the only way to get a number divisible by $3$ by adding three of these is by adding the three ones. So those must go first. Now we have an odd sum, and since the next average must be divisible by $4, 71$ must be next. That leaves $80$ for last, so the answer is $\mathrm{C}$.

 



\textbf{Solution 3}

we know that the average of the scores is an integer 

 so that means s1+s2+s3+s4 must be an even number divisible by 4 

 we have 3 even scores and 2 odd scores

 which means that the last score cannot be odd because otherwise, we would get an odd number divided by an even number in the denominator.

 so we have answers that are even 

 76,80,82

 We see 3 cases where 76 is the last score, 80 is the last score, and 82 is the last score  76= 1 mod(5)which means 80+82+71+91= 0 mod(4)80+82+71+91= 4 mod(5)

 80= 0 mod(5)76+82+71+91= 0 mod(5)76+82+71+91= 0 mod(4)

 82= 2 mod(5)76+80+71+91= 0 mod(4)76+80+71+91= 3 mod(5)

 Case 1: 76324 is divisible by 4 324 divided by 5 is 1which means 76 is not the last number

 case 2:324-4=320320 is divisible by 4 320 is divisible by 5 which means this case is true. 

 case 3:320-2=318318 is not divisible by 4 which makes it incorrect even though 

 it has a remainder of 3 when divided by 5

 



\textbf{Solution 4}

The test numbers we have include 71, 76, 80, 82, 91 The answers are same numbers: 71, 76, 80, 82, 91 

 You could either subtract everything by 66 to make things easier5,10,14,16,25getting 70 when adding up now add 5,10,13,16,25 and it easier to see which one is divisible by 6

 or you could  perform the following action test number minus 6[test number/6]

 5,4,2,4,1add them up to get 6+6+4now we have to add another 2 for it to be divisible by 6

 before the operation 2 was actually 80 thus 80 is the answer

 


	\item (\textit{AMC 10 10 2000}) Two non-zero real numbers, $a$ and $b$, satisfy $ab=a-b$.  Find a possible value of $\frac{a}{b}+\frac{b}{a}-ab$.


 $\textbf{(A)}\ -2 \qquad\textbf{(B)}\ -\frac{1}{2} \qquad\textbf{(C)}\ \frac{1}{3} \qquad\textbf{(D)}\ \frac{1}{2} \qquad\textbf{(E)}\ 2$


 
 \textbf{Answer: }E



\textbf{Solution 1}

$\frac {a}{b} + \frac {b}{a} - ab = \frac{a^2 + b^2}{ab} - (a - b) = \frac{a^2 + b^2}{a-b} - \frac{(a-b)^2}{(a-b)} = \frac{2ab}{a-b} = \frac{2(a-b)}{a-b} =\boxed{\textbf{(E)}\ 2}$. 

 



\textbf{Solution 2}

Adding $b$ to both sides of the equation relating $a$ and $b$ yields $ab + b = a$, so $b(a+1) = a$ and therefore $b = \frac{a}{a+1}$.  

 Substituting $b=\frac{a}{a+1}$ into the desired expression yields $\frac{a}{\frac{a}{a+1}} + \frac{\frac{a}{a+1}}{a} - a \cdot{\frac{a}{a+1}}$ $= a+1+\frac{1}{a+1}-\frac{a^2}{a+1} = a+1+\frac{1-a^2}{a+1} = a+1+\frac{(1-a)(1+a)}{1+a}$ $=a+1+1-a=\boxed{\textbf{(E)}\ 2}$ 

 



\textbf{Solution 3}

This simplifies to $ab+b-a=0 \Rightarrow (a+1)(b-1) = -1$. The two integer solutions to this are $(-2,2)$ and $(0,0)$. The problem states than $a$ and $b$ are non-zero, so we consider the case of $(-2,2)$. So, we end up with $\frac{-2}{2} + \frac{2}{-2} - 2 \cdot -2 = 4 - 1 - 1 = 2 \Rightarrow \boxed{\text{E}}$

 



\textbf{Solution 4}

Just realize that two such numbers are $a = 1$ and $b = \frac{1}{2}$. You can see this by plugging in $a = 1$ and then solving for b. With this, you can solve and get $2 \Rightarrow\boxed{\text{E}}$

 



\textbf{Solution 5}

Set $a$ to some nonzero number. In this case, I'll set it to $4$. 

 Then solve for $b$. In this case, $b=0.8$. 

 Now just simply evaluate. In this case it's 2. So since 2 is a possible value of the original expression, select $\boxed{E}$.

 



\textbf{Solution 6}

Notice that $a=\frac{a}{b}-1$ and $b=1-\frac{b}{a}$. Then, $\frac{a}{b}=1+a$ and $\frac{b}{a}=1-b$.

 $\frac{a}{b}+\frac{b}{a}-ab=(1+a)+(1-b)-(a-b)=2$. The answer is $\boxed{E}$.

 



\textbf{Solution 7}

We are told $ab=a-b$.

 Dividing by $a$, we get $b = 1 - \frac{b}{a} \Rightarrow \boxed{ \frac{b}{a} = 1 - b}$.

 We can also solve for $\frac{a}{b}$ by dividing by $b$. Dividing, we get $a = \frac{a}{b} - 1 \Rightarrow \boxed{ \frac{a}{b} = a + 1}$. 

 Plugging these values into $\frac{a}{b} + \frac{b}{a} - ab$, we get:

 
 \[(a + 1) + (1 - b) - (a - b) = 2\]

 Therefore, the final answer is $\boxed{\text{E}}$

 -Dustin\_Links

 


	\item (\textit{AMC 10 10 2000}) The diagram shows $28$ lattice points, each one unit from its nearest neighbors. Segment $AB$ meets segment $CD$ at $E$. Find the length of segment $AE$.


 
                    \begin{figure}[H]
                        \centering
                        \includegraphics[width=0.5\textwidth]{Problem/AMC_10_10_2000_16_Problem_Figure_1.png}
                    \end{figure}
                    


 $\textbf{(A)}\ \frac{4\sqrt{5}}{3} \qquad\textbf{(B)}\ \frac{5\sqrt{5}}{3} \qquad\textbf{(C)}\ \frac{12\sqrt{5}}{7} \qquad\textbf{(D)}\ 2\sqrt{5} \qquad\textbf{(E)}\ \frac{5\sqrt{65}}{9}$


 
 \textbf{Answer: }B



\textbf{Solution 1}

Let $l_1$ be the line containing $A$ and $B$ and let $l_2$ be the line containing $C$ and $D$. If we set the bottom left point at $(0,0)$, then $A=(0,3)$, $B=(6,0)$, $C=(4,2)$, and $D=(2,0)$. 

 The line $l_1$ is given by the equation $y=m_1x+b_1$. The $y$-intercept is $A=(0,3)$, so $b_1=3$. We are given two points on $l_1$, hence we can compute the slope, $m_1$ to be $\frac{0-3}{6-0}=-\frac{1}{2}$, so $l_1$ is the line $y=\frac{-1}{2}x+3$

 Similarly, $l_2$ is given by $y=m_2x+b_2$. The slope in this case is $\frac{2-0}{4-2}=1$, so $y=x+b_2$. Plugging in the point $(2,0)$ gives us $b_2=-2$, so $l_2$ is the line $y=x-2$.

 At $E$, the intersection point, both of the equations must be true, so 
 \begin{align*} y=x-2, y=\frac{-1}{2}x+3 &\Rightarrow x-2=\frac{-1}{2}x+3 \\ &\Rightarrow x=\frac{10}{3} \\ &\Rightarrow y=\frac{4}{3} \\ \end{align*}


 We have the coordinates of $A$ and $E$, so we can use the distance formula here: 
 \[\sqrt{\left(\frac{10}{3}-0\right)^2+\left(\frac{4}{3}-3\right)^2}=\frac{5\sqrt{5}}{3}\]

 which is answer choice $\boxed{\text{B}}$

 



\textbf{Solution 2}


                    \begin{figure}[H]
                        \centering
                        \includegraphics[width=0.5\textwidth]{Solution/AMC_10_10_2000_16_Solution_Figure_1.png}
                    \end{figure}
                    

 Draw the perpendiculars from $A$ and $B$ to $CD$, respectively. As it turns out, $BC \perp CD$. Let $F$ be the point on $CD$ for which $AF\perp CD$.

 $m\angle AFE=m\angle BCE=90^\circ$, and $m\angle AEF=m\angle CEB$, so by AA similarity, 
 \[\triangle AFE\sim \triangle BCE \Rightarrow \frac{AF}{AE}=\frac{BC}{BE}\]

 By the Pythagorean Theorem, we have $AB=\sqrt{3^2+6^2}=3\sqrt{5}$, $AF=\sqrt{2.5^2+2.5^2}=2.5\sqrt{2}$, and $BC=\sqrt{2^2+2^2}=2\sqrt{2}$. Let $AE=x$, so $BE=3\sqrt{5}-x$, then 
 \[\frac{2.5\sqrt{2}}{x}=\frac{2\sqrt{2}}{3\sqrt{5}-x}\] 
 \[x=\frac{5\sqrt{5}}{3}\]

 This is answer choice $\boxed{\text{B}}$

 



\textbf{Solution 3}


                    \begin{figure}[H]
                        \centering
                        \includegraphics[width=0.5\textwidth]{Solution/AMC_10_10_2000_16_Solution_Figure_2.png}
                    \end{figure}
                    

 Drawing line $\overline{BD}$ and parallel line $\overline{CF}$, we see that $\triangle FCE \sim \triangle BDE$ by AA similarity. Thus $\frac{FE}{EB} = \frac{FC}{DB} = \frac{2}{4} = \frac{1}{2}$. Reciprocating, we know that $\frac{EB}{FE} = 2$ so $\frac{EB+FE}{FE} = 2+1 \Rightarrow \frac{FB}{FE} = 3$. Reciprocating again, we have $\frac{FE}{FB} = \frac{1}{3} \Rightarrow FE = \frac{1}{3}FB$. We know that $FD = 2$, so by the Pythagorean Theorem, $FB = \sqrt{2^{2} + 4^{2}} = 2\sqrt{5}$. Thus $FE = \frac{1}{3}FB = \frac{2\sqrt{5}}{3}$. Applying the Pythagorean Theorem again, we have $AF = \sqrt{1^{2}+2^{2}} = \sqrt{5}$. We finally have $AE = AF + FE = \sqrt{5} + \frac{2\sqrt{5}}{3} = \frac{5\sqrt{5}}{3} \Rightarrow \boxed{\text{B}}$

 



\textbf{Solution 4}


                    \begin{figure}[H]
                        \centering
                        \includegraphics[width=0.5\textwidth]{Solution/AMC_10_10_2000_16_Solution_Figure_3.png}
                    \end{figure}
                    

 Extend line $\overline{DC}$ as above.  This creates two similar triangles whose side lengths have the ratio $5:4$.  Therefore $AE=\frac{5}{9}AB$.  Using Pythagorean theorem to find $AB$ gives us:

 
 \[AE=\frac{5}{9}AB=\frac{5}{9}\sqrt{3^2+6^2}=\frac{5}{9}\sqrt{45}= \boxed{\textbf{(B) }\frac{5\sqrt{5}}{3}}\]

 


	\item (\textit{AMC 10 10 2000}) Boris has an incredible coin changing machine. When he puts in a quarter, it returns five nickels; when he puts in a nickel, it returns five pennies; and when he puts in a penny, it returns five quarters. Boris starts with just one penny. Which of the following amounts could Boris have after using the machine repeatedly?


 $\textbf{(A) } \$3.63  \qquad \textbf{(B) } \$5.13 \qquad \textbf{(C) } \$6.30 \qquad \textbf{(D) } \$7.45 \qquad \textbf{(E) } \$9.07$


 
 \textbf{Answer: }D



\textbf{Solution}

If Boris puts a quarter or nickel in the machine, he neither gains nor loses any money.  If he puts a penny in the machine, he gains $\$1.24$.

 Let $x$ be the number of times Boris puts a penny in the machine.  Then he has $1.24x+0.01$ dollars.  Since $x$ is an integer, the only answer choice Boris could have is $\boxed{\textbf{(D)}\ \$7.45}$, which happens if $x = 6$.  Boris could achieve this by putting a penny, a quarter, a nickel, and then five pennies into the machine.

 


	\item (\textit{AMC 10 10 2000}) Charlyn walks completely around the boundary of a square whose sides are each $5$ km long. From any point on her path she can see exactly $1$ km horizontally in all directions. What is the area of the region consisting of all points Charlyn can see during her walk, expressed in square kilometers and rounded to the nearest whole number?


 $\textbf{(A) } 24 \qquad\textbf{(B)}\ 27 \qquad\textbf{(C)}\ 39 \qquad\textbf{(D)}\ 40 \qquad\textbf{(E)}\ 42$


 
 \textbf{Answer: }C



\textbf{Solution}

The area she sees looks at follows:

 
                    \begin{figure}[H]
                        \centering
                        \includegraphics[width=0.5\textwidth]{Solution/AMC_10_10_2000_18_Solution_Figure_1.png}
                    \end{figure}
                    

 The part inside the walk has area $5\cdot 5 - 3\cdot 3 = 16$. The part outside the walk consists of four rectangles, and four arcs. Each of the rectangles has area $5\cdot 1=5$. The four arcs together form a circle with radius $1$. 

 Therefore the total area she can see is $16 + 4\cdot 5 + \pi\cdot 1^2 = 36+\pi \simeq 39.14$, which rounded to the nearest integer is $39$. $\boxed{C}$

 



\textbf{Solution 2(Common sense)}

The area of the bigger square when drawing out is basically 49, 7\textasciicircum{}2. The smaller area that Charlyn cannot see it 9, 3\textasciicircum{}2. So 49-9=40, but we have to remember that the area of the bigger rounded square isn't exactly 49(rounded edges), so the closest answer choice that's less than 40 is $39$ $\boxed{C}$~dnathan12

 


	\item (\textit{AMC 10 10 2000}) Through a point on the hypotenuse of a right triangle, lines are drawn parallel to the legs of the triangle so that the triangle is divided into a square and two smaller right triangles. The area of one of the two small right triangles is $m$ times the area of the square. The ratio of the area of the other small right triangle to the area of the square is


 $\textbf{(A)}\ \frac{1}{2m+1} \qquad\textbf{(B)}\ m \qquad\textbf{(C)}\ 1-m \qquad\textbf{(D)}\ \frac{1}{4m} \qquad\textbf{(E)}\ \frac{1}{8m^2}$


 
 \textbf{Answer: }D



\textbf{Solution 1}


                    \begin{figure}[H]
                        \centering
                        \includegraphics[width=0.5\textwidth]{Solution/AMC_10_10_2000_19_Solution_Figure_1.png}
                    \end{figure}
                    

 WLOG, let a side of the square be $1$. Simple angle chasing shows that the two right triangles are \href{https://artofproblemsolving.com/wiki/index.php?title=Similar\_triangles}{similar}. Thus the ratio of the sides of the triangles are the same. Since $A = \frac{1}{2}bh = \frac{h}{2}$, the base of the triangle with area $m$ is $2m$. Therefore $\frac{2m}{1} = \frac{1}{x}$ where $x$ is the height of the other triangle. $x = \frac{1}{2m}$, and the area of that triangle is $\frac{1}{2} \cdot 1 \cdot \frac{1}{2m} = \frac{1}{4m}\ \text{\boxed{D}}$.

 


	\item (\textit{AMC 10 10 2000}) Let $A$, $M$, and $C$ be nonnegative integers such that $A+M+C=10$.  What is the maximum value of $A\cdot M\cdot C+A\cdot M+M\cdot C+C\cdot A$?


 $\textbf{(A)}\ 49 \qquad\textbf{(B)}\ 59 \qquad\textbf{(C)}\ 69 \qquad\textbf{(D)}\ 79 \qquad\textbf{(E)}\ 89$


 
 \textbf{Answer: }C



\textbf{Solution 1}

The trick is to realize that the sum $AMC+AM+MC+CA$ is similar to the product $(A+1)(M+1)(C+1)$. If we multiply $(A+1)(M+1)(C+1)$, we get 
 \[(A+1)(M+1)(C+1) = AMC + AM + AC + MC + A + M + C + 1.\] We know that $A+M+C=10$, therefore $(A+1)(M+1)(C+1) = (AMC + AM + MC + CA) + 11$ and 
 \[AMC + AM + MC + CA = (A+1)(M+1)(C+1) - 11.\] Now consider the maximal value of this expression. Suppose that some two of $A$, $M$, and $C$ differ by at least $2$.  Then this triple $(A,M,C)$ is not optimal. (To see this, WLOG let $A\geq C+2.$ We can then increase the value of $(A+1)(M+1)(C+1)$ by changing $A \to A-1$ and $C \to C+1$. Or, we can consider the equality case of the AM-GM inequality, which would maximize this product; intuitively, the closest we can get to that equality case is what we have below.)

 Therefore the maximum is achieved when $(A,M,C)$ is a permutation of $(3,3,4)$. The value of $(A+1)(M+1)(C+1)$ in this case is $4\cdot 4\cdot 5=80,$ and thus the maximum of $AMC + AM + MC + CA$ is $80-11 = \boxed{\textbf{(C)}\ 69}.$

 



\textbf{Solution 2}

Notice that if we want to maximize $AMC + AM + MC + AC$, we want A, M, and C to be as close as possible. For example, if $A = 7, B = 2,$ and $C=1,$ then the expression would have a much smaller value than if we were to substitute $A = 4, B = 5$, and $C = 1$. So to make A, B, and C as close together as possible, we divide $\frac{10}{3}$ to get $3$. Therefore, A must be 3, M must be 3, and C must be 4.  $AMC + AM + MC + AC = 36 + 12 + 12 + 9 = 69$. So the answer is $\boxed{\textbf{(C)}\ 69.}$

 



\textbf{Solution 3}

According to our knowledge in AM-GM, the closer numbers are, the value of their product is bigger. Assume that $A=B<C$, we can get the set $A=3,B=3,C=4$ which the answer is $\boxed{\textbf{(C)}\ 69.}$~bluesoul

 



\textbf{Solution 4 (Lagrange Multipliers / Calculus)}

This problem can be modeled as an optimization problem and solved using multivariable calculus.

 Our objective function is defined as:
 \[f(A, M, C) = AMC + AM + MC + AC\]

 Subject to the constraint:
 \[g(A, M, C) = A + M + C - 10 = 0\]

 We set up the Lagrangian function $\mathcal{L}(A, M, C, \lambda)$ by subtracting the product of the multiplier $\lambda$ and the constraint function from the objective function:

 
 \[\mathcal{L}(A, M, C, \lambda) = f(A, M, C) - \lambda \cdot g(A, M, C)\]
 \[\mathcal{L}(A, M, C, \lambda) = AMC + AM + MC + AC - \lambda(A + M + C - 10)\]
 \[\mathcal{L}(A, M, C, \lambda) = AMC + AM + MC + AC - \lambda A - \lambda M - \lambda C + 10\lambda\]

 The critical points of this system are found by setting the gradient $\nabla\mathcal{L}$ equal to the zero vector. To compute the partial derivatives, we differentiate with respect to one variable while treating all other variables as constants. Setting each partial derivative to $0$ gives us the following system of equations:

 
 \[\frac{\partial \mathcal{L}}{\partial A} = MC + M + C - \lambda = 0 \implies \lambda = MC + M + C \quad \text{(Eq. 1)}\]
 \[\frac{\partial \mathcal{L}}{\partial M} = AC + A + C - \lambda = 0 \implies \lambda = AC + A + C \quad \text{(Eq. 2)}\]
 \[\frac{\partial \mathcal{L}}{\partial C} = AM + A + M - \lambda = 0 \implies \lambda = AM + A + M \quad \text{(Eq. 3)}\]
 \[\frac{\partial \mathcal{L}}{\partial \lambda} = -(A + M + C - 10) = 0 \implies A + M + C = 10 \quad \text{(Eq. 4)}\]

 By equating $\text{Eq. 1}$ and $\text{Eq. 2}$, we get:
 \[MC + M + C = AC + A + C\]
 \[MC + M = AC + A\]
 \[MC + M - AC - A = 0\]

 Grouping terms and factoring gives:
 \[(M - A)(C + 1) = 0\]

 This equation holds if $M = A$ or $C = -1$. Since $A, M,$ and $C$ are given to be non-negative integers, $C$ cannot be $-1$, meaning only the former condition ($M = A$) can hold.

 By symmetry, equating $\text{Eq. 2}$ and $\text{Eq. 3}$ yields $A = C$. Thus, the critical point occurs when all three variables are equal:
 \[A = M = C\]

 Substituting this into the original constraint ($\text{Eq. 4}$):
 \[3A = 10 \implies A = M = C = \frac{10}{3}\]

 In the real domain, the function reaches its absolute maximum at $\left(\frac{10}{3}, \frac{10}{3}, \frac{10}{3}\right)$. Since $A, M,$ and $C$ must be integers, we look for the closest points of integral coordinates that satisfy the constraint. This is achieved by distributing the values as evenly as possible: $3, 3,$ and $4$ (in any order).

 Plugging these integer values into our objective function:
 \[f(3, 3, 4) = (3)(3)(4) + (3)(3) + (3)(4) + (3)(4) = 36 + 9 + 12 + 12 = 69\]

 Thus, the maximum value is $\boxed{\textbf{(C)}\ 69.}$.

 


	\item (\textit{AMC 10 10 2000}) If all alligators are ferocious creatures and some creepy crawlers are alligators, which statement(s) must be true?


 
 \[\textrm{I. All alligators are creepy crawlers.}\]

 \[\textrm{II. Some ferocious creatures are creepy crawlers.}\]

 \[\textrm{III. Some alligators are not creepy crawlers.}\]


 $\textbf{(A)}\ \text{I only} \qquad\textbf{(B)}\ \text{II only} \qquad\textbf{(C)}\ \text{III only} \qquad\textbf{(D)}\ \text{II and III only} \qquad\textbf{(E)}\ \text{None must be true}$


 
 \textbf{Answer: }B



\textbf{Solution}

Statement I could be false.  A possible situation in which statement I is false is if the only creatures are two alligators, one of which is a creepy crawler and one of which is not.

 Statement II must be true.  Since some creepy crawlers are alligators and all alligators are ferocious creatures, it follows that some creepy crawlers are ferocious creatures, and thus some ferocious creatures are creepy crawlers.

 Statement III could be false.  A possible situation in which statement III is false is if the only creature is a creepy crawler which is also an alligator.

 Therefore, only $\boxed{\textbf{(B)}\text{ II only}}$ is necessarily true.

 


	\item (\textit{AMC 10 10 2000}) One morning each member of Angela's family drank an 8-ounce mixture of coffee with milk. The amounts of coffee and milk varied from cup to cup, but were never zero. Angela drank a quarter of the total amount of milk and a sixth of the total amount of coffee. How many people are in the family?


 $\textbf{(A)}\ 3 \qquad\textbf{(B)}\ 4 \qquad\textbf{(C)}\ 5 \qquad\textbf{(D)}\ 6 \qquad\textbf{(E)}\ 7$


 
 \textbf{Answer: }C



\textbf{Solution 1}

If Angela drank $\frac{1}{4}$ of the milk and $\frac{1}{4}$ of the coffee, she would have drunk $\frac{1}{4}$ of the ingredients.  However, she actually drank $\frac{1}{6}$ of the coffee, so she drank less than $\frac{1}{4}$ of the ingredients.

 Similarly, if she drank $\frac{1}{6}$ of the milk and $\frac{1}{6}$ of the coffee, she would have drunk $\frac{1}{6}$ of the ingredients.  However, she actually drank $\frac{1}{4}$ of the milk, so she drank more than $\frac{1}{6}$ of the ingredients.

 If there are $n$ people in the family, then each person drank $\frac{1}{n}$ of the ingredients.  So $\frac{1}{6} < \frac{1}{n} < \frac{1}{4}$, and thus $4 < n < 6$.  Since $n$ is an integer, $n = \boxed{\textbf{(C)}\ 5}$.

 



\textbf{Solution 2}

Let $c$ be the total amount of coffee, $m$ of milk, and $p$ the number of people in the family. Then each person drinks the same total amount of coffee and milk (8 ounces), so 
 \[\left(\frac{c}{6} + \frac{m}{4}\right)p = c + m\]Regrouping, we get $2c(6-p)=3m(p-4)$. Since both $c,m$ are positive, it follows that $6-p$ and $p-4$ are also positive, which is only possible when $p = 5\ \mathrm{(C)}$.

 



\textbf{Solution 3}

Again, let $c,$ $m,$ and $p$ be the total amount of coffee, total amount of milk, and number of people in the family, respectively. Then the total amount that is drunk is $8p,$ and also $c+m.$ Thus, $c+m = 8p,$ so $m = 8p-c$ and $c = 8p-m.$

 We also know that the amount Angela drank, which is $\frac{c}{6} + \frac{m}{4},$ is equal to $8$ ounces, thus $\frac{c}{6} + \frac{m}{4} = 8.$ Rearranging gives  
 \[24p - c = 96.\] Now notice that $c > 0$ (by the problem statement). In addition, $m > 0,$ so $c = 8p-m < 8p.$ Therefore, $0 < c < 8p,$ and so $24p > 24p-c > 16p.$ We know that $24p-c = 96,$ so 
 \[24p > 96 > 16p.\] From the leftmost inequality, we get $p > 4,$ and from the rightmost inequality, we get $p < 6.$ The only possible value of $p$ is $p = 5\ \mathrm{(C)}$.

 



\textbf{Solution 4}

Let $c,$ $m,$ and $p$ be the total amount of coffee, total amount of milk, and number of people in the family, respectively. $c$ and $m$ obviously can't be $0$. We know $\frac{c}{6} + \frac{m}{4} = 8$ or $2c + 3m = 96$ and $c + m = 8p$ or $2c + 2m = 16p$. Then, 
 \[(2c + 2m) + m = 16p + m = 96\]Because $16p$ and $96$ are both divisible by $16$, $m$ must also be divisible by $16$. Let $m = 16k$. Now, 
 \[3(16k) + 2c = 48k + 2c = 96\] $k$ can't be $0$, otherwise $m$ is $0$, and $k$ can't be $2$, otherwise $c$ is $0$. Therefore $k$ must be $1$, $m = 16$ and $c = 24$. $c + m = 24 + 16 = 40 = 8p$. Therefore, $p = 5\ \mathrm{(C)}$.

 



\textbf{Solution 5}

Let $m$, $c$ be the total amounts of milk and coffee, respectively. In order to know the number of people, we first need to find the total amount of mixture $t = m + c$. We are given that 
 \[\frac{m}{4} + \frac{c}{6} = 8\]Multiplying the equation by $4$ yields 
 \[m + \frac{2}{3}c = (m + c) - \frac{1}{3}c = t - \frac{1}{3}c = 32\]Since $\frac{1}{3}c > 0$, we have $t > 32$. Now multiplying the equation by $6$ yields 
 \[\frac{3}{2}m + c = (m+c) + \frac{1}{2}m = t + \frac{1}{2}m = 48\]Since $\frac{1}{2}m > 0$, we have $t < 48$. Thus, $32 < t < 48$.

 Since $t$ is a multiple of $8$, the only possible value for $t$ in that range is $40$. Therefore, there are $\frac{40}{8} = 5$ people in Angela's family. $\mathrm{(C)}$.

 



\textbf{Solution 6 (Constraints)}

If there were 4 people in the family, and each of them drank exactly the same amount of coffee and milk as Angela, there would be too much coffee. If there were 6 people in the family, and each of them drank exactly the same amount of coffee and milk as Angela, there would be not enough milk. Thus, it has to be $\boxed{5}$.

 (But how can we assume that each of them drank exactly the same amount of coffee and milk?)

 



\textbf{Solution 7 (Guess and check)}

The number of ounces that Angela's family drank has to be a multiple of 8, so we can find the right answer by guessing random values for the number of ounces of coffee and milk Angela drank. With Angela drinking 4 ounces of milk and 4 ounces of coffee, we get 40 total ounces Angela's family drank. Dividing that by 8(each person drank 8 ounces) we get 5 members who drank the coffee with milk, or $\boxed{C}$  ~ Mathyguy88

 


	\item (\textit{AMC 10 10 2000}) When the mean, median, and mode of the list 
 \[10,2,5,2,4,2,x\] are arranged in increasing order, they form a non-constant arithmetic progression. What is the sum of all possible real values of $x$?


 $\textbf{(A)}\ 3 \qquad\textbf{(B)}\ 6 \qquad\textbf{(C)}\ 9 \qquad\textbf{(D)}\ 17 \qquad\textbf{(E)}\ 20$


 



 
 \textbf{Answer: }E



\textbf{Solution 1}


                \begin{itemize}
                    \item The mean is $\frac{10+2+5+2+4+2+x}{7} = \frac{25+x}{7}$.\item Arranged in increasing order, the list is $2,2,2,4,5,10$, so the median is either $2,4$ or $x$ depending upon the value of $x$.\item The mode is $2$, since it appears three times.
                \end{itemize}


                We apply casework upon the median:

 
                \begin{itemize}
                    \item If the median is $2$ ($x \le 2$), then the arithmetic progression must be constant.\item If the median is $4$ ($x \ge 4$), because the mode is $2$, the mean can either be $0,3,6$ to form an arithmetic progression. Solving for $x$ yields $-25,-4,17$ respectively, of which only $17$ works because it is larger than $4$.\item If the median is $x$ ($2 \le x \le 4$), we must have the arithmetic progression $2, x, \frac{25+x}{7}$.  Thus, we find that $2x=2+\frac{25+x}{7}$ so $x=3$.
                \end{itemize}


                The answer is $3 + 17 = \boxed{20\ \mathrm{(E)}}$.-Error fixed(BIGFOOT09) Changes Copyrighted©

 



\textbf{Solution 2}

The mode will always be $2$, as there are already $3$ of them.  The median will be $2$, $3$, or $4$, as only $1$ new number is added.  The mean will be all the numbers added together divided by $7$.

 However:The median can't actually be $2$, as that makes it no longer an arithmetic progression.  Therefore, the median has to be $3$ or $4$.

 When the median is $3$, $x$ has to equal $3$ as well, as it becomes the median.  Therefore, the mean has to be $4$, which also happens to work when $x=3$.

 When the median is $4$, $x$ has to be greater than $3$, as well as allowing the mean to equal $6$.  As the current values added together equal $25$, we need a number that will allow us to make the mean equal $6 \cdot 7 = 42$.  As $42-25=17$, $x=17$.

 The question asks for the sum of all solutions for $x$, so the answer is $3 + 17 = \boxed{20\ \mathrm{(E)}}$.~Particle Physics

 


	\item (\textit{AMC 10 10 2000}) Let $f$ be a function for which $f\left(\frac{x}{3}\right)=x^2+x+1$. Find the sum of all values of $z$ for which $f(3z)=7$.


 $\textbf{(A)}\ -\frac{1}{3} \qquad\textbf{(B)}\ -\frac{1}{9} \qquad\textbf{(C)}\ 0 \qquad\textbf{(D)}\ \frac{5}{9} \qquad\textbf{(E)}\ \frac{5}{3}$


 
 \textbf{Answer: }B



\textbf{Solution 1}

Let $y = \frac{x}{3}$; then $f(y) = (3y)^2 + 3y + 1 = 9y^2 + 3y+1$. Thus $f(3z)-7=81z^2+9z-6=3(9z-2)(3z+1)=0$, and $z = -\frac{1}{3}, \frac{2}{9}$. These sum up to $\boxed{\textbf{(B) }-\frac19}$.

 



\textbf{Solution 2 (Similar)}

This is quite trivially solved, as $3x = \dfrac{9x}{3}$, so $P(3x) = P(9x/3) = 81x^2 + 9x + 1 = 7$. $81x^2+9x-6 = 0$ has solutions $-\frac{1}{3}$ and $\frac{2}{9}$. Adding these yields a solution of $\boxed{\textbf{(B) }-\frac19}$.

 



\textbf{Solution 3}

Similar to Solution 1, we have $=81z^2+9z-6=0.$ The answer is the sum of the roots, which by \href{https://artofproblemsolving.com/wiki/index.php?title=Vieta\%27s\_Formulas}{Vieta's Formulas} is  $-\frac{b}{a}=-\frac{9}{81}=\boxed{\textbf{(B) }-\frac19}$.

 



\textbf{Solution 4}

Set $f\left(\frac{x}{3} \right) = x^2+x+1=7$ to get $x^2+x-6=0.$ From either finding the roots (-3 and 2), or using Vieta's formulas, we find the sum of these roots to be $-1.$ Each root of this equation is $9$ times greater than a corresponding root of $f(3z) = 7$ (because $\frac{x}{3} = 3z$ gives $x = 9z$), thus the sum of the roots in the equation $f(3z)=7$ is $-\frac{1}{9}$ or $\boxed{\textbf{(B) }-\frac19}$.

 



\textbf{Solution 5}

Since we have $f(x/3)$, $f(3z)$ occurs at $x=9z.$ Thus, $f(9z/3) = f(3z) = (9z)^2 + 9z + 1$. We set this equal to 7:

 $81z^2 + 9z +1 = 7 \Longrightarrow 81z^2 + 9z - 6 = 0$. For any quadratic $ax^2 + bx +c = 0$, the sum of the roots is $-\frac{b}{a}$. Thus, the sum of the roots of this equation is $-\frac{9}{81} = \boxed{\textbf{(B) }-\frac19}$.

 


	\item (\textit{AMC 10 10 2000}) In year $N$, the $300^\text{th}$ day of the year is a Tuesday. In year $N+1$, the $200^\text{th}$ day is also a Tuesday. On what day of the week did the $100^\text{th}$ day of year $N-1$ occur?


 $\textbf{(A)}\ \text{Thursday} \qquad\textbf{(B)}\ \text{Friday} \qquad\textbf{(C)}\ \text{Saturday} \qquad\textbf{(D)}\ \text{Sunday} \qquad\textbf{(E)}\ \text{Monday}$


 
 \textbf{Answer: }A



\textbf{Solution 1}

There are either 
 \[65 + 200 = 265\] or 
 \[66 + 200 = 266\] days between the first two dates depending upon whether or not year $N+1$ is a leap year (since the February 29th of the leap year would come between the 300th day of year $N$ and 200th day of year $N + 1$). Since $7$ divides into $266$ but not $265$, for both days to be a Tuesday, year $N$ must be a leap year.

 Hence, year $N-1$ is not a leap year, and so since there are 
 \[265 + 300 = 565\] days between the date in years $N,\text{ }N-1$, this leaves a remainder of $5$ upon division by $7$. Since we are subtracting days, we count 5 days before Tuesday, which gives us  $\boxed{\mathbf{(A)} \ \text{Thursday}.}$

 



\textbf{Solution 2}

The $300-29\cdot 7=97^{\text{th}}$ day of year $N$ and the $200-15\cdot 7=95^{\text{th}}$ day of year $N+1$ are Tuesdays. If there were the same number of days in year $N$ and year $N-1,$ the $99^{\text{th}}$ day of year $N-1$ will be a Tuesday. But year $N$ is a leap year because $97-95 \cong 366 \pmod{7},$ so the $98^{\text{th}}$ day of year $N-1$ is a Tuesday. It follows that the $100^{\text{th}}$ day of year $N-1$ is $\boxed{\mathbf{(A)} \ \text{Thursday}.}$

 


\end{enumerate}